\chapter{Generalized Priming}
\label{chap:priming}

\scaffold{Chapter intro: the second contribution. \revizor{}'s priming step cannot
report a whole class of genuine leaks---those whose only effect on the run that
causes them lies outside the L1D cache. We give a formal model, prove the limitation,
show the generalization keeps the violation definition and detects strictly more, and
give two optimal localizers. Draw from the priming-extension note and uasc2027
``Generalized priming''.}

\section{The detection gap}
\label{sec:priming-gap}
\scaffold{Intuition: an input can write a microarchitectural resource (a BTB way, a
TLB entry, the prefetcher, a predictor) that its own run never reads back into the
cache. The write leaves no mark on the run's own L1D trace and survives only in the
carried-over state, surfacing in a later run. Priming was built to cancel
carried-over state, so it cancels the evidence together with the noise. Cross-run
Spectre-v2 and the canonicality leak are members of this class.}

\section{A formal model of priming}
\label{sec:priming-model}
\scaffold{Definitions: run and batch, the initialization sequence, the contract
trace and the two hardware traces $\HL$ (L1D) and $\HF$ (full state), determinism, the
violation of a contract, the current priming decision rule, input classes, and
boosting.}

\section{Lanes and splices}
\label{sec:priming-lanes}
\scaffold{The lane / lane-boosting matrix (rows are lanes/batches, columns are
classes), the prefix/suffix splice, and the splice measurement $r(t)$.}

\section{Current priming misses noncache violations}
\label{sec:priming-limitation}
\scaffold{State and prove the theorem that the current rule never reports a noncache
violation---using only the shape of the rule, not any particular channel. Then show
the canonicality leak and cross-run Spectre-v2 are noncache violations.}

\section{Noncache differences reach the cache}
\label{sec:priming-observe}
\scaffold{Every observable state difference is exposed in the L1D by some generated
test case (register-dependent load), so the generalization loses no observable
difference---it reads the leaked residue through a later run's $\HL$, never $\HF$.}

\section{Localizing the leak}
\label{sec:priming-localize}
\scaffold{A boundary exists when the end readings differ; every boundary is a genuine
violation (two inputs of one class separated by one identical probe); self-dependent
versus cross-input leaks; running both directions.}

\section{Two localizers}
\label{sec:priming-localizers}
\scaffold{Linear scan (returns the largest leaking class, worst case $k$ runs) and
binary search (returns one leaking class, $\lceil\log_2(k+1)\rceil$ runs); state
optimality, defer full proofs and the worked binary-search figure to
\autoref{appendix:localization}.}
